\documentclass{article}
\usepackage{amsmath,amssymb,amsthm}
\newtheorem{theorem}{Theorem}
\newtheorem{lemma}{Lemma}
\newcommand{\CE}{\operatorname{CE}}
\newcommand{\I}{\operatorname{I}}
\begin{document}
\title{An output-only reference repairing the repeated-label comb transcript}
\author{A011, B03; stage 20261011T232009Z}
\maketitle
All logarithms are to base two unless written $\ln$; cubes are uniform.
The result repairs an existing complete-code cost, not universal FEI.

\begin{lemma}[Binary empty-sector correction]
For $p+q=1$, $p,q\ge0$, with $L(t)=t\log(1/t)$ and $L(0)=0$,
\[
 L(p^2)+L(q^2)\le4pq.
\]
\end{lemma}
\begin{proof}
By symmetry assume $0\le p\le1/2$. Set
$A(p)=p^2\log(1/p)+(1-p)^2\log(1/(1-p))$ and
$E(p)=2p(1-p)-A(p)$. Then $E(0)=E(1/2)=0$ and
\[
 E'(0)=2-1/\ln2>0,\quad E'(1/2)=0,
 \quad E''(p)=-4+\frac{2\ln(p(1-p))+6}{\ln2},
 \quad E'''(p)=\frac{2}{\ln2}\left(\frac1p-\frac1{1-p}\right)>0.
\]
Also $E''(0)=-\infty$ and
$E''(1/2)=(6-8\ln2)/\ln2>0$. Thus $E''$ changes sign once,
$E'$ first decreases then increases to zero, and $E$ first increases then
decreases to zero. Consequently $E\ge0$. Since
$L(p^2)+L(q^2)=2A(p)$, the result follows. The elementary bounds
$1/2<\ln2<3/4$ justify the stated signs: the lower bound follows by
integrating $1/t$ on $[1,2]$, and the upper bound by the strict trapezoid
bound for the convex function $1/t$ on that interval.
\end{proof}

\begin{theorem}[All repeated two-label combs]
Let $a\ge1$, and let $\sigma_1,\ldots,\sigma_{a+1}\in\{-1,1\}$
be arbitrary. Query independent selectors $x_1,\ldots,x_a$ in order,
stopping at the first positive coordinate; index $a+1$ is the terminal
all-negative leaf. A positive label selects $\chi_A(y)$, and a negative
label selects $\chi_B(y)$, for distinct nonempty sets $A,B$ on the same
independent data cube. The data supports can overlap. Call the result $F$.

There is an explicit normalized output spectral reference $Q_F$, fixed
from the labelled tree, whose full cost satisfies
\[
 \CE(\widehat F^2,Q_F)=H(F)
 \le C_*\I_{\rm sel}(F),\qquad
 C_*=4+\frac{2}{e\ln2}<\frac{41}{8}.
\]
It admits a physical output prefix code of expected length at most
$(41/8)\I(F)$. No boundary-retention hypothesis is needed, including when
$\I_{\rm sel}/D\to0$, where $D=2(1-2^{-a})$.
\end{theorem}
\begin{proof}
Define physical suffix means and disagreement probabilities by
\[
 m_{a+1}=\sigma_{a+1},\qquad
 m_j=(\sigma_j+m_{j+1})/2,\qquad
 r_j=(1-\sigma_jm_{j+1})/2,
 \qquad p=(1+m_1)/2,\quad q=1-p.
\]
For the Boolean selector-label sign $g$, repeated restriction gives
\[
 \widehat g(T)=2^{1-j}\sigma_jr_j(-1)^{|T|-1},
 \qquad T\ne\varnothing,\quad j=\max T.
\]
Indeed, all earlier selectors continue on the negative branch, producing
factors $(1-x_i)/2$; at query $j$ the coefficient of $x_j$ is
$(\sigma_j-m_{j+1})/2=\sigma_jr_j$. This proves the formula directly.
Now $F=(1+g)\chi_A/2+(1-g)\chi_B/2$. Its empty-selector masses are
$p^2,q^2$. For every $j$, its $2^j$ max-index-$j$ output atoms each have
mass $2^{-2j}r_j^2$. Their total is
\[
 v_j=2^{-j}r_j^2,
 \qquad p^2+q^2+\sum_j v_j=1.
\]
The latter follows from the explicit Fourier expansion and Parseval of
the actual Boolean $F$, including zero atoms. A normalized generator
selects either empty atom or a group $j$ with these masses. Conditional
on $j$, choose a fair data label and a uniform subset of $[j-1]$, then
include $j$. This is $Q_F$, with no physical stopping label. All its
parameters were computed from the physical suffix means.

Write $J=\I_{\rm sel}(F)$. The physical boundary calculation and the
uniform predecessor-subset calculation give respectively
\[
 J=\sum_j2^{-j}r_j,
 \qquad J=\frac12\sum_j(j+1)v_j.
\]
For the first, reaching query $j$ has probability $2^{-(j-1)}$, its
continuation has a different character with probability $r_j$, and
different unsigned characters disagree on half the data cube. For the
second, each group has expected selector cardinality $(j+1)/2$.
Put $V=\sum_jv_j=2pq$. The complete output cost is
\[
 L(p^2)+L(q^2)+2\sum_j jv_j
       +2\sum_{j:r_j>0}v_j\log(1/r_j).
\]
The lemma bounds its empty part by $2V$, while its middle part is exactly
$4J-2V$. Thus no empty-sector or normalization term is left unpaid.
The scalar maximum $r\log(1/r)\le1/(e\ln2)$ gives
\[
 2\sum_j2^{-j}r_j^2\log(1/r_j)
 \le\frac{2}{e\ln2}\sum_j2^{-j}r_j
 =\frac{2}{e\ln2}J.
\]
This proves the ideal reference cost bound. The rational constant follows
from $e>8/3$ (the first four terms of its power series sum to $8/3$ and
the remaining terms are positive) and $\ln2>2/3$ (strict midpoint
integral bound for the strictly convex function $1/t$ on $[1,2]$).

Assign lengths $\lceil\log(1/Q_F(S))\rceil$ to the finite positive
outputs. Their Kraft sum is at most one, so a canonical prefix trie
provides a decoder. If the law is deterministic use the empty word.
Otherwise the expected output length is less than $H(F)+1$. Since
$\I_{\rm data}(F)=p|A|+q|B|\ge1$, this is at most
$(41/8)\I(F)$. A prefix code of the \emph{output} is asserted, not a
shortened code decoding the old physical witness.
\end{proof}

\paragraph{Failure of passive marginalization.}
Set all stopping labels to $\chi_A$ and only the terminal to $\chi_B$.
Let $u=2^{-a}$. The old path-then-uniform-mask Kraft reference has
\[
 Q_{\rm raw}(A)=\sum_{j=1}^a4^{-j}=(1-u^2)/3.
\]
The true output mass there is $(1-u)^2$. Consequently
\[
 \CE(P_F,Q_{\rm raw})\ge(1-u)^2\log\frac3{1-u^2}\longrightarrow\log3,
 \qquad \I_{\rm sel}(F)=au\longrightarrow0.
\]
Marginalizing the old uniform-Kraft reference is insufficient. The repaired
reference uses the coherent suffix calculation. It does not subtract
posterior entropy from a transmitted physical-leaf code.

\paragraph{Independent reconstruction of the A012 objection.}
For the different Boolean function $F_r=1+A_r(y-1)$, $u=2^{-r}$,
the control-first tree has $D=2(1-u)$, deterministic spectral leaves,
and raw cost $2D$. Direct physical edge counting gives
$\I(F_r)=u(r+1)$. Hence raw cost/influence is
$4(2^r-1)/(r+1)\to\infty$. Its Fourier constant coefficient is $1-u$
and every other occupied coefficient has magnitude $u$, giving
$H(F_r)=2ru(2-u)+L((1-u)^2)\le4ru+2u/\ln2$.
If the exact witness retains the physical leaf $L$, then
\[
 H(W\mid S)\ge H(L\mid S)\ge H(L)-H(S)
 \ge2(1-u)-4ru-2u/\ln2\longrightarrow2.
\]
This accepts only the specified transcript-obstruction subclaim of
A012 v1, not its entire hereditary/class package. The comb theorem
concerns two nonempty unsigned data labels; it is not a proved code
for every constant/Boolean-core activation. A012 owns that separate repair.

\paragraph{Arbitrary seed and the universal line.}
For any Boolean $h$ on disjoint coordinates and any complete seed reference
$Q_h$, the product $Q_F\otimes Q_h$ is a complete subprobability and
\[
 \CE(P_{Fh},Q_F\otimes Q_h)=\CE(P_F,Q_F)+\CE(P_h,Q_h),\qquad
 \I(Fh)=\I(F)+\I(h).
\]
For $h^{\otimes t}$ the latter terms multiply by $t$; no seed entropy is
erased. More directly, for an arbitrary Boolean selector $g$ the exact
two-character ledger is
\[
 H(F_g)=L(p^2)+L(q^2)+\operatorname{Var}(g)+\tfrac12 H_{\ne\varnothing}(g),
 \qquad \I_{\rm sel}(F_g)=\I(g)/2.
\]
For balanced $g$ this is $H(F_g)=2+H(g)/2$. The first unsupported
universal line is a cheap full conditional reference replacing the uniform
predecessor law in the comb's max-index groups for arbitrary Boolean
selectors. Its conditional entropy is not paid by merely compressing a
physical label. Full FEI, literature novelty and independent acceptance
of the present repair are not claimed.
\end{document}
